% Options for packages loaded elsewhere
\PassOptionsToPackage{unicode}{hyperref}
\PassOptionsToPackage{hyphens}{url}
\documentclass[
  11pt,
  a4paper,
]{article}
\usepackage{xcolor}
\usepackage[margin=1in]{geometry}
\usepackage{amsmath,amssymb}
\setcounter{secnumdepth}{-\maxdimen} % remove section numbering
\usepackage{iftex}
\ifPDFTeX
  \usepackage[T1]{fontenc}
  \usepackage[utf8]{inputenc}
  \usepackage{textcomp} % provide euro and other symbols
\else % if luatex or xetex
  \usepackage{unicode-math} % this also loads fontspec
  \defaultfontfeatures{Scale=MatchLowercase}
  \defaultfontfeatures[\rmfamily]{Ligatures=TeX,Scale=1}
\fi
\usepackage{lmodern}
\ifPDFTeX\else
  % xetex/luatex font selection
  \setmainfont[]{STIXGeneral}
  \setmonofont[]{Menlo}
\fi
% Use upquote if available, for straight quotes in verbatim environments
\IfFileExists{upquote.sty}{\usepackage{upquote}}{}
\IfFileExists{microtype.sty}{% use microtype if available
  \usepackage[]{microtype}
  \UseMicrotypeSet[protrusion]{basicmath} % disable protrusion for tt fonts
}{}
\makeatletter
\@ifundefined{KOMAClassName}{% if non-KOMA class
  \IfFileExists{parskip.sty}{%
    \usepackage{parskip}
  }{% else
    \setlength{\parindent}{0pt}
    \setlength{\parskip}{6pt plus 2pt minus 1pt}}
}{% if KOMA class
  \KOMAoptions{parskip=half}}
\makeatother
\usepackage{longtable,booktabs,array}
\newcounter{none} % for unnumbered tables
\usepackage{calc} % for calculating minipage widths
% Correct order of tables after \paragraph or \subparagraph
\usepackage{etoolbox}
\makeatletter
\patchcmd\longtable{\par}{\if@noskipsec\mbox{}\fi\par}{}{}
\makeatother
% Allow footnotes in longtable head/foot
\IfFileExists{footnotehyper.sty}{\usepackage{footnotehyper}}{\usepackage{footnote}}
\makesavenoteenv{longtable}
\setlength{\emergencystretch}{3em} % prevent overfull lines
\providecommand{\tightlist}{%
  \setlength{\itemsep}{0pt}\setlength{\parskip}{0pt}}
\usepackage{newunicodechar}
\newfontfamily\numeralfallback{STIX Two Text}
\newunicodechar{∎}{\ensuremath{\blacksquare}}
\usepackage{fvextra}
\DefineVerbatimEnvironment{Highlighting}{Verbatim}{breaklines,breakanywhere,breakanywheresymbolpre={},breakanywheresymbolpost={},breaksymbolleft={},breaksymbolright={},commandchars=\\\{\}}
\newunicodechar{⁰}{{\numeralfallback\char"2070\relax}}
\newunicodechar{₀}{{\numeralfallback\char"2080\relax}}
\newunicodechar{₁}{{\numeralfallback\char"2081\relax}}
\newunicodechar{₂}{{\numeralfallback\char"2082\relax}}
\newunicodechar{₃}{{\numeralfallback\char"2083\relax}}
\newunicodechar{₄}{{\numeralfallback\char"2084\relax}}

\usepackage{seqsplit}
\renewcommand{\texttt}[1]{{\ttfamily\seqsplit{#1}}}
\setlength{\emergencystretch}{6em}

\newunicodechar{⁻}{{\numeralfallback\char"207B\relax}}
\newunicodechar{⁸}{{\numeralfallback\char"2078\relax}}
\newunicodechar{ⁱ}{{\numeralfallback\char"2071\relax}}
\newunicodechar{⁴}{{\numeralfallback\char"2074\relax}}
\newunicodechar{⁷}{{\numeralfallback\char"2077\relax}}
\newunicodechar{⁶}{{\numeralfallback\char"2076\relax}}
\newunicodechar{⁵}{{\numeralfallback\char"2075\relax}}
\newunicodechar{₈}{{\numeralfallback\char"2088\relax}}
\newunicodechar{ᵢ}{\textsubscript{i}}
\newunicodechar{ⱼ}{\textsubscript{j}}
\newunicodechar{ₜ}{\textsubscript{t}}
\usepackage{bookmark}
\IfFileExists{xurl.sty}{\usepackage{xurl}}{} % add URL line breaks if available
\urlstyle{same}
% fallback for those not using the hyperref driver hyperxmp:
\makeatletter
\@ifundefined{xmpquote}{\newcommand{\xmpquote}[1]{#1}}{}
\makeatother
\hypersetup{
  hidelinks,
  pdfcreator={LaTeX via pandoc}}

\author{}
\date{}

\begin{document}

\section{Equivariant incidence reconstruction for two Livingstone-graph
polytopes}\label{equivariant-incidence-reconstruction-for-two-livingstone-graph-polytopes}

\subsection{1. Scope and source
boundary}\label{scope-and-source-boundary}

Hartley, Hubard and Leemans {[}HHL12{]} construct two regular rank-four
polytopes from the Livingstone graph of the first Janko group. Their
main contribution is geometric: the faces are explicit graph subsets,
and incidence is specified by intersections and distances. The existence
of abstract polytopes with the same automorphism group does not
establish this construction. The additional task is to identify the
literal graph incidence relations with the already certified coset
incidence relations.

We give a common verification argument and carry out that
identification. All six relations defining the second structure, of type
\(\{5,6,5\}\), agree with its coset geometry. For the first structure,
of type \(\{5,3,5\}\), five of the six displayed relations agree. The
remaining printed relation omits a compatibility condition: it creates
twelve incident facets at each 2-face, whereas the intended construction
has two. We exhibit a pairwise incident flag that cannot extend to a
chamber. Adding the compatibility supplied by the source's Proposition
3.5 recovers the intended regular polytope.

Consequently this work does \textbf{not} prove both main theorems for
the literal structures printed in {[}HHL12{]}. It recovers Theorem 5.1
and its stabilizer lattice, proves a corrected form of Theorem 4.1, and
records a defect in the printed definition. The correction must remain
visible in any source ledger or compression count. It is not an
open-conjecture result. No priority for the correction is asserted: a
targeted title/DOI search on 10 October 2026 did not locate an erratum,
but that is not a proof that none exists.

The reusable mathematical step is an elementary equivariant
reconstruction. The finite source-specific part is a complete
examination of the graph objects and their neighbour sets. The abstract
marked-group certificates of Module C are retained unchanged; they
supply generation, the full intersection property, typed Helly and the
two Schläfli types. Named-group data, exact permutation-group algorithms
and the classical identification of the degree-266 action remain
explicit inputs.

\subsection{2. Reconstruction from one
neighbourhood}\label{reconstruction-from-one-neighbourhood}

Let \(G\) act on the right on nonempty sets \(X_i\), transitively on
each type. For distinct types, let \(R_{ij}\subseteq X_i\times X_j\) be
a \(G\)-invariant binary relation. Choose objects \(x_i\in X_i\), put
\(H_i=G_{x_i}\), and suppose \(x_i R_{ij}x_j\) for every pair of
distinct types.

\textbf{Theorem 2.1 (equivariant reconstruction).} Suppose that, for
every \(i\ne j\), \[
 \{y\in X_j:x_iR_{ij}y\}=x_j^{H_i}.
 \tag{2.1}
\] Then the maps \[
 H_i g\longmapsto x_i^g
\] give an isomorphism from the right-coset incidence structure of the
subgroups \(H_i\) to the given typed incidence structure. If the
neighbour sets are finite, it is enough to check \[
 \#\{y\in X_j:x_iR_{ij}y\}
   =[H_i:H_i\cap H_j].
 \tag{2.2}
\]

\emph{Proof.} Equality of right cosets is equivalent to equality of the
associated images of \(x_i\), because \(H_i\) is exactly its stabilizer.
Transitivity gives surjectivity. Two cosets \(H_i g,H_j h\) intersect if
and only if there is \(a\in G\) with \(x_i^a=x_i^g\) and
\(x_j^a=x_j^h\). Thus coset incidence is exactly the diagonal orbit of
\((x_i,x_j)\).

The base pair lies in \(R_{ij}\), so invariance puts that entire
diagonal orbit inside \(R_{ij}\). Conversely, suppose \(x_i^gR_{ij}y\).
Applying \(g^{-1}\) gives \(x_iR_{ij}y^{g^{-1}}\). By (2.1),
\(y^{g^{-1}}=x_j^u\) for some \(u\in H_i\). Then
\((x_i^g,y)=(x_i^{ug},x_j^{ug})\), so the pair lies in the diagonal
orbit. This proves equality of the binary relations. All flags and
residues are defined by these relations, and therefore are preserved as
well.

For the finite criterion, invariance of \(R_{ij}\) and the base
incidence give \(x_j^{H_i}\subseteq R_{ij}(x_i)\). Orbit--stabilizer
gives \(|x_j^{H_i}|=[H_i:H_i\cap H_j]\). Equality of the two finite
cardinalities proves equality of the sets. \(\square\)

The hypothesis in (2.2) concerns the \emph{complete} neighbour set.
Counting a chosen orbit of neighbours would merely repeat
orbit--stabilizer and would prove nothing about possible additional
incidences. The omitted relation in {[}HHL12{]} gives a concrete example
of this distinction.

The three Lean declarations \texttt{reconstruction},
\texttt{localOrbit\_subset} and \texttt{reconstruction\_of\_card} in
\texttt{phase179\_formal/EquivariantIncidence.lean} formalize the
diagonal-orbit argument and the finite-cardinality step, with left
actions. Passing to inverse coordinates gives the right-action
convention above. They do not formalize the named graph, the whole coset
construction, the J1 certificates, or either complete source theorem.

\subsection{3. The graph and its finite object
families}\label{the-graph-and-its-finite-object-families}

Let \(G=J_1\) in its standard transitive action on 266 points. We retain
the classical identification of this action, with vertex stabilizer
\(\operatorname{PSL}_2(11)\), and of its unique orbital graph of valency
11 with the Livingstone graph \(L\). This is the construction used in
{[}HHL12, §3{]}, ultimately from Livingstone and the classical
distance-regular-graph literature. The finite checker reads the named
permutation group from TomLib and takes the unique 11-point orbit of a
point stabilizer as the neighbourhood of that point. Transport defines
the other neighbourhoods.

The checker verifies graph invariance under every supplied group
generator, symmetry, absence of loops, and the complete intersection
array \[
 \{11,10,6,1;1,1,5,11\}.
\] It checks the intersection numbers for every ordered pair of
vertices. In particular the distance distribution is
\(1,11,110,132,12\), and the graph has girth five. There are 1463 edges.
No source geometry list or source pentagon list is used as input.

Every pentagon is enumerated by taking its least vertex first and its
smaller adjacent vertex second, extending four edges, and requiring the
last edge to close. Girth five excludes chords; the ordering removes the
two orientations and the five starting points. The enumeration produces
11,704 pentagons, and their actual \(G\)-orbits have lengths 2926 and
8778. Following {[}HHL12{]}, call the smaller orbit white,
\(\mathcal W\), and the larger orange, \(\mathcal O\).

For an edge \(e\), its stabilizer has vertex orbits of lengths \[
 2,10,20,24,30,60,120.
\] Let \(P_e\) be its unique orbit of length 10. The induced subgraph is
checked to have strongly regular parameters \((10,3,0,1)\), and hence is
a Petersen graph. Transporting one edge and its 10-set gives 1463
distinct sets \(P_e\). Each contains six white pentagons; their
complements in \(P_e\) are orange pentagons. Each white pentagon belongs
to three of these Petersen graphs.

\textbf{Lemma 3.1 (completeness of the Petersen family).} Every induced
Petersen subgraph of \(L\) is one of the \(P_e\).

\emph{Proof.} Every Petersen graph contains a pentagon
\(p_0p_1p_2p_3p_4\). Its five remaining vertices \(q_i\) are joined to
\(p_i\) by the five spokes, and their cyclic order is
\(q_0,q_2,q_4,q_1,q_3\). Choose \(q_0\) among the neighbours of \(p_0\)
outside the pentagon. There are nine choices. Thereafter \(q_2\) must be
a common neighbour of \(q_0,p_2\), and is unique if it exists: two
distinct common neighbours would give a 4-cycle. The same argument
determines \(q_4,q_1,q_3\) in turn. Reject a repeated vertex or a
failure to close the final edge. This exhausts every possible Petersen
extension of the chosen pentagon.

Apply this nine-choice procedure to one white and one orange pentagon.
It gives respectively three and one extensions, all in the constructed
family. Every pentagon is in one of these two \(G\)-orbits, and the
family \(\{P_e\}\) is invariant. Therefore every extension of every
pentagon lies in this family. \(\square\)

The construction of the Petersen family has a useful source correction.
The literal statement of {[}HHL12, Proposition 3.3{]} describes its ten
vertices as all vertices at distance two from both endpoints of \(e\).
That common distance-two set has \textbf{40} vertices. For adjacent
\(u,v\), the intersection array gives this number without enumeration.
For each \(z\in\Gamma_2(u)\), exactly four neighbours of \(u\) are at
distance two from \(z\), because \(u\in\Gamma_2(z)\) and \(a_2=4\).
Summing over the 110 choices of \(z\) and using transitivity on the
eleven neighbours of \(u\) gives \(110\cdot4/11=40\). The checker
independently finds 40.

The preceding paragraph in {[}HHL12, §3.1.4{]} states the needed
additional condition. The correct 10-set is \[
 P_e=\{z:d(z,u)=d(z,v)=2,
       \text{ some }W\in\mathcal W\text{ contains }e\cup\{z\}\}.
 \tag{3.1}
\] We verify (3.1) for one edge. Both sides are equivariant, and \(G\)
is edge-transitive, so this proves it for every edge. We do not
attribute this condition to a new idea: it is already present in the
source paragraph.

For each \(W\subset P_e\), write \(O=P_e\setminus W\). The other two
Petersen graphs containing \(W\) are \(P_{e_0},P_{e_1}\). Define the
source quadruple \[
 q=(W,O,e,\{e_0,e_1\}).
\] The family \(\mathcal Q\) contains 8778 quadruples and is one
\(G\)-orbit. The program constructs it from all 1463 Petersen sets and
their six white pentagons, and separately checks transitivity. It does
not assume that identical edge data identifies a quadruple: there are
six different quadruples over each edge.

The final family uses the other conjugacy class of \(A_5\) in a vertex
stabilizer. The complete maximal-subgroup classes of
\(G_v\cong\operatorname{PSL}_2(11)\) contain two \(A_5\) classes, whose
neighbour orbit lengths are respectively \(1,10\) and \(5,6\). The
second gives eleven six-subsets of \(N(v)\); they form a
\(2\!-(11,6,3)\) design. Write \[
 t=(v,V_5,V_6),\qquad V_5=N(v)\setminus V_6.
\] Their full \(G\)-orbit \(\mathcal T\) has size 2926. The checker
verifies every pair of the eleven neighbours lies in exactly three
blocks. It also verifies {[}HHL12, Proposition 3.4{]}: a triple of
neighbours lies in one or two blocks, and lies in two precisely when the
triple together with \(v\) is contained in a Petersen graph. This uses
all 165 neighbour triples and the complete Petersen family proved above.

For \(u\in V_6\), put \(e=\{u,v\}\). The intersection of the stabilizer
of \(t\) with the pointwise edge stabilizer has order 10. Its two orbits
on \(P_e\) have length five, one white and one orange, and determine
exactly one quadruple over \(e\). The checker verifies all six choices
of \(u\) for one triple; transitivity covers all triples. This recovers
Proposition 3.5 and defines the \textbf{induced quadruple} of \((t,e)\).

There is an entirely graph-based description of the same compatibility.
For \(z\in P_e\), let \(\phi_v(z)\) be the unique common neighbour of
\(v,z\). Equation (3.1) and girth five show this is the neighbour
appearing in the white pentagon containing \(e,z\). The map
\(\phi_v:P_e\to N(v)\setminus\{u\}\) is the equivariant correspondence
used in the source. The induced-quadruple condition is \[
 \{\phi_v(W),\phi_v(O)\}
    =\{V_5,V_6\setminus\{u\}\}.
 \tag{3.2}
\] The unordered pair allows either pentagon colour to correspond to the
five-set. No choice of colour is silently imposed.

\subsection{4. The literal definition and its
repair}\label{the-literal-definition-and-its-repair}

The source defines the first type sets as \[
 X_0=E(L),\quad X_1=X_2=\mathcal Q,
 \quad X_3=\mathcal T,
\] where an edge in \(X_0\) stands for its pair \((e,P_e)\). The second
construction has the same first three type sets and \(X_3=V(L)\). The
definitions of \(\mathcal Q,\mathcal T\) above are exactly those in the
source; they carry no implicit incidence condition from another type.

For the first construction, {[}HHL12, p.205{]} declares
\(q=(W',O',e',\{e'_0,e'_1\})\) incident to \(t=(v,V_5,V_6)\) when \[
 v\in e',\qquad |V_6\cap e'|=1.
 \tag{4.1}
\] Every point of the local \(2\!-(11,6,3)\) design belongs to six
blocks. Thus, for either endpoint \(v\) of \(e'\), there are six
admissible choices of \(V_6\), namely all blocks containing the other
endpoint. Equation (4.1) does not involve \(W'\) or \(O'\). It therefore
gives \textbf{twelve}, not two, facets at every 2-face.

This contradicts the conclusion of the source's subsequent calculation
on p.208. That paragraph uses the correspondence of §3.2.4 to select a
block from the \emph{whole quadruple}, while the displayed relation uses
only its edge. Proposition 3.5 gives a quadruple from a block and edge;
it does not make a block unique when only the edge is prescribed. The
distinction is essential because six quadruples share that edge.

The exact group check decomposes the twelve literal neighbours into
stabilizer orbits of lengths 2 and 10. More strongly, it supplies the
following explicit extra flag, in the degree-266 TomLib labelling: \[
\begin{aligned}
 q={}&(\{16,63,157,214,249\},\{67,88,143,165,230\},
          \{42,135\},\{\{5,238\},\{86,191\}\}),\\
 t={}&(42,\{8,47,102,111,237\},
          \{71,78,128,135,226,234\}).
\end{aligned}
\] The checker verifies that these are actual members of \(\mathcal Q\)
and \(\mathcal T\), respectively, and satisfy (4.1). It evaluates the
other two relevant \emph{literal} source predicates against every one of
the 8778 possible type-one faces. None is incident to both \(q\) and
\(t\). Hence this pairwise incident flag cannot extend to a chamber. The
literal incidence structure is not a geometry in the source's own sense.

Define the corrected first structure \(\widehat Q\) by retaining the
five other source relations and replacing (4.1) by (4.1)
\textbf{together with} the induced-quadruple condition (3.2). This is an
explicit mathematical repair, not a reinterpretation claimed to be
present in the printed formula. The second structure \(P\) is left
literally unchanged.

\subsection{5. Identification with the retained marked
geometries}\label{identification-with-the-retained-marked-geometries}

The Module C certificate \texttt{phase40\_compress\_operations\_check.g}
constructs a marking \(s=(a,b,c,d)\) in the same named group, with pair
orders \[
 (|ab|,|ac|,|ad|,|bc|,|bd|,|cd|)=(5,2,2,3,2,5).
\] Writing \(w=(bcd)^5\), its second marking is \((a,wb,c,d)\), with
pair orders \((5,2,2,6,2,5)\). These are the two previously certified
cases of types \(\{5,3,5\}\) and \(\{5,6,5\}\). That certificate proves
generation, all 256 subgroup-meet identities, the typed coset Helly
conditions, trivial chamber stabilizer and order-two panel stabilizers.
Its proof and finite input are retained, rather than being inferred from
the graph count.

For either marking let \(H_i\) be the subgroup generated by all markers
except the \(i\)-th. The new checker finds the source face fixed by each
\(H_i\); there is exactly one candidate of each type. It then calculates
the full stabilizer of that object and checks equality with \(H_i\). In
particular it does not infer equality merely from the fact that the
subgroup fixes the object.

Each selected base pair satisfies the appropriate source relation, using
the repaired relation only for the \(\{2,3\}\) pair of \(\widehat Q\).
For every pair \(i<j\), the checker evaluates the predicate on
\textbf{every member of \(X_j\)} with the chosen \(x_i\) fixed. It
compares this entire set with the \(H_i\)-orbit of \(x_j\), not just
their sizes. The outcomes are:

Each entry below gives the neighbour count followed by the pair
stabilizer order.

{\def\LTcaptype{none} % do not increment counter
\begin{longtable}[]{@{}lrr@{}}
\toprule\noalign{}
Types \(i,j\) & \(\widehat Q\) & \(P\) \\
\midrule\noalign{}
\endhead
\bottomrule\noalign{}
\endlastfoot
0,1 & 12 / 10 & 12 / 10 \\
0,2 & 30 / 4 & 30 / 4 \\
0,3 & 20 / 6 & 10 / 12 \\
1,2 & 5 / 4 & 5 / 4 \\
1,3 & 5 / 4 & 5 / 4 \\
2,3 & 2 / 10 & 2 / 10 \\
\end{longtable}
}

All predicates are invariant because they use graph distances, graph
subsets, membership, intersections and the equivariant correspondence
(3.2). Theorem 2.1 identifies each entire structure with its retained
coset geometry.

\textbf{Theorem 5.1 (source recovery with explicit correction).} The
literal second construction \(P\) of {[}HHL12{]} is a regular rank-four
polytope of type \(\{5,6,5\}\), with face counts
\((1463,8778,8778,266)\). The corrected first construction
\(\widehat Q\) is a regular rank-four polytope of type \(\{5,3,5\}\),
with face counts \((1463,8778,8778,2926)\). Each has full
type-preserving automorphism group \(J_1\). Their flag stabilizer
lattices are respectively those printed in source Figures 10 and 9.

\emph{Proof.} The exact stabilizer and neighbour-set checks above give
Theorem 2.1's hypotheses. The retained C certificate and recognition
theorem therefore prove flag transitivity, thinness and residual
connectedness. The marker product orders give the indicated polygon
sections, so the string geometries are the stated abstract regular
polytopes. Orbit--stabilizer gives the displayed face counts.

The four parabolic orders are \((120,20,20,60)\) for \(\widehat Q\) and
\((120,20,20,660)\) for \(P\), with types
\(2\times A_5,D_{20},D_{20},A_5\) or \(\operatorname{PSL}_2(11)\). In
the order of the table, the pair stabilizers are
\(D_{10},C_2^2,S_3,C_2^2,C_2^2,D_{10}\) for \(\widehat Q\); only \(S_3\)
changes, to \(C_2\times S_3\cong D_{12}\), for \(P\). Each triple
stabilizer is \(C_2\), and the fourfold stabilizer is trivial. The
checker prints all sixteen intersections, including the empty
intersection \(G\). These are exactly the two printed stabilizer
lattices.

Finally, a type-preserving automorphism fixing a chamber fixes its
adjacent chamber across every panel, because each panel belongs to
exactly two chambers. Chamber connectedness propagates this to every
chamber, and every face belongs to a chamber. Thus such an automorphism
is the identity. The full type-preserving automorphism group acts freely
on chambers and has order at most the number of chambers, namely
\(|J_1|\). The already faithful, regular action of \(J_1\) reaches this
bound. \(\square\)

This last argument uses the geometry's thinness and connectedness,
rather than assuming that every automorphism of the geometry comes from
the graph. It therefore supplies the full-group conclusion without an
unproved converse between graph and geometry automorphisms.

\subsection{6. Attribution, verification and
cost}\label{attribution-verification-and-cost}

The constructions, their intended polytope types and their geometric
ingredients belong to {[}HHL12{]}. Its Propositions 3.1--3.3 explicitly
cite the classical distance-regular-graph literature. The \(J_1\)
action, vertex subgroup structure, local design and previous
abstract-polytope classification are retained historical inputs.
Aschbacher's theorem appears as source Theorem 2.1; our argument instead
reuses the existing U/C recognition proof. The orbit--stabilizer
reconstruction in §2 is standard. The present work contributes its
common verification interface, the exact source-to-coset reconstruction,
and the explicit audit and repair of the missing condition. None of
these statements asserts historical priority.

The computational proof obligations are finite and fully specified:
construct every relevant family; prove its completeness and transitivity
by the arguments above; calculate the exact stabilizers; compare every
local source neighbour set with the appropriate stabilizer orbit. The
GAP checker performs those tasks from the named group, not from the
source geometry lists. It also retains the literal failure as an
expected counterexample, so a successful repaired run cannot erase it.

The Lean theorem certifies the reusable logical implication after those
hypotheses have been supplied. It is not a Lean certificate of the GAP
objects or the two full polytopes. The trusted base includes the pinned
compiler and cached Mathlib dependency objects for the component proof,
and GAP/TomLib and their named-group data for the finite work. The
classical graph identification and the retained C proof remain
mathematical inputs.

This application must be charged for the new manuscript, GAP program,
exact drivers and Lean source, together with the retained C/U
dependencies under the project's common allocation rule. Source PDFs,
generated renderings, logs, manifests and review reports are support
artifacts. The source has 22 pages; that number alone does not determine
a compression ratio. In particular the literal first principal theorem
cannot be admitted as an unchanged corollary when its displayed
definition is false. The source ledger records one recovered principal
construction and one explicitly repaired construction, without inflating
the unchanged principal-source count.

\subsection{References}\label{references}

{[}HHL12{]} M. I. Hartley, I. Hubard and D. Leemans, \emph{Regular
4-polytopes from the Livingstone graph of Janko's first group}, Journal
of Algebraic Combinatorics \textbf{35} (2012), 193--214. DOI
\href{https://doi.org/10.1007/s10801-011-0300-x}{10.1007/s10801-011-0300-x}.
Exact primary edition: \texttt{sources/catalogue\_L068.pdf}, SHA256
\texttt{e40b3ee73dec3754248a7c9b659fc3015f8ee59c2b98ce46141baae8fe825ad9}.
An accessible primary mirror is \href{https://emis.de/ft/43964}{EMIS}.

{[}C{]} Retained Module C source, exact program and review bindings for
the J1 marking and twist constructions; the phase179 manifest identifies
the required retained files. The complete source subgroup lattice is not
re-proved or newly counted here.

\end{document}
